\chapter{Contextul Problemei și Cerințele Domeniului}

\section{Problematica deșeurilor în spații verzi și publice}

Spațiile verzi, parcurile, aleile pietonale și zonele publice deschise fac parte din infrastructura de bază a unui oraș. Ele nu sunt doar zone decorative, ci locuri folosite constant de cetățeni: pentru deplasare, odihnă, activități sportive sau interacțiune socială. Calitatea acestor spații influențează direct imaginea orașului și confortul celor care le utilizează. Din această perspectivă, prezența deșeurilor abandonate reprezintă o problemă practică, vizibilă și greu de ignorat.

În mod obișnuit, deșeurile întâlnite în astfel de zone sunt obiecte mici sau medii, precum ambalaje, hârtii, pungi, doze metalice, sticle ori resturi alimentare. Dimensiunea redusă face ca ele să fie ușor de aruncat, dar și ușor de pierdut în fundalul scenei. O pungă poate avea o formă neregulată, o hârtie poate fi confundată cu o zonă luminată de pe alee, iar o sticlă transparentă poate deveni greu vizibilă în funcție de unghiul camerei. Aceste detalii, aparent simple pentru un observator aflat la fața locului, devin mai dificile atunci când scena este analizată automat dintr-un flux video.

Problema nu este doar una de curățenie. Din punct de vedere administrativ, administratorii spațiilor publice trebuie să identifice zonele în care apar frecvent incidente, să intervină rapid și să poată documenta situațiile relevante. În lipsa unei monitorizări eficiente, activitatea de întreținere rămâne reactivă: se intervine după ce problema este observată, raportată sau acumulată. În schimb, un sistem capabil să semnaleze rapid apariția unui posibil incident poate ajuta la reducerea timpului de reacție și la organizarea mai bună a intervențiilor.

În această lucrare, termenul de deșeu desemnează un obiect vizibil în imagine, aparținând unei categorii materiale relevante pentru monitorizare: plastic, hârtie, sticlă, metal sau alte materiale similare. Lucrarea nu urmărește întregul proces de management al deșeurilor, cum ar fi colectarea, transportul, reciclarea sau planificarea infrastructurii de salubritate. Accentul este pus pe o etapă mai restrânsă, dar importantă: identificarea automată a obiectelor de tip deșeu în cadre video și semnalarea situațiilor în care apariția lor poate indica un incident de aruncare sau abandonare.

Această delimitare este esențială: problema abordată nu este doar detecția statică a deșeurilor, ci interpretarea lor în context video. Sistemul trebuie să observe modificarea scenei, să coreleze apariția obiectului cu prezența unei persoane și să ofere utilizatorului responsabil o dovadă locală care poate fi verificată.

Din perspectiva viziunii artificiale, recunoașterea deșeurilor în contexte naturale este o problemă dificilă. Setul de date TACO evidențiază diversitatea mare a obiectelor de tip litter și dificultățile produse de iluminare, fundal, perspectivă și ocluzie \cite{proenca2020taco}. Studii recente dedicate detecției deșeurilor în medii naturale și urbane confirmă faptul că problema nu se reduce la clasificarea unor obiecte izolate, ci presupune date realiste, categorii materiale variate și evaluare în contexte vizuale dificile \cite{majchrowska2022waste}. Spre deosebire de seturi generale precum COCO sau PASCAL VOC, unde multe categorii de obiecte au forme relativ stabile și apar frecvent în prim-plan \cite{lin2014microsoft}, deșeurile din spații publice pot fi mici, deformate, parțial acoperite sau foarte asemănătoare cu elementele mediului. Acest lucru justifică folosirea unei soluții adaptate domeniului și testate pe scenarii apropiate de utilizarea dorită.

\section{Necesitatea monitorizării automate prin analiză video}

Monitorizarea clasică a spațiilor verzi se realizează prin patrule, verificări periodice, sesizări ale cetățenilor sau camere video. Fiecare dintre aceste metode are utilitate, dar niciuna nu rezolvă complet problema observării continue. Patrulele sunt limitate de timp și personal, sesizările apar după producerea incidentului, iar camerele video devin eficiente doar dacă există un mod practic de analiză a materialului înregistrat. O cameră care doar stochează imagini nu reduce, în sine, efortul de verificare.

Analiza video automată schimbă rolul camerei. Aceasta nu mai este doar o sursă de înregistrare, ci o sursă de date care poate fi prelucrată în timp aproape real. Fiecare cadru poate fi analizat pentru localizarea obiectelor relevante, iar detecțiile pot fi urmărite în timp pentru a observa dacă scena s-a modificat. În loc ca utilizatorul responsabil să urmărească permanent fluxul, sistemul poate semnala doar momentele care merită verificate și poate salva dovezile asociate.

Totuși, automatizarea nu elimină rolul omului. În contextul lucrării de față, sistemul nu este proiectat pentru a lua decizii juridice și nu stabilește sancțiuni. El funcționează ca un instrument de asistență: detectează, filtrează, salvează dovezi și propune incidente pentru validare. Utilizatorul responsabil verifică imaginea sau clipul asociat și decide dacă incidentul este confirmat, respins ca fals pozitiv sau arhivat. Această separare este importantă deoarece modelele de detecție pot greși, mai ales în scene reale, unde iluminarea, mișcarea, ocluziile și fundalul sunt greu de controlat.

Pentru ca o astfel de soluție să fie utilă, ea trebuie să îndeplinească mai multe cerințe. Sistemul trebuie să proceseze cadre video într-un timp rezonabil, să ofere rezultate ușor de verificat, să păstreze dovezi locale și să permită administrarea clară a incidentelor. În același timp, trebuie să trateze cu atenție datele video, deoarece acestea pot conține persoane și informații sensibile. În lucrare, aceste cerințe sunt reflectate prin folosirea unui pipeline de detecție și clasificare, a unei logici temporale pentru semnalarea incidentelor și a unei aplicații web locale pentru monitorizare, validare și administrare.

Aceste cerințe arată că problema nu poate fi redusă la o singură etapă de inferență. Detecția obiectelor este necesară, dar nu suficientă. Sistemul trebuie să includă și mecanisme de interpretare temporală, interfață de validare, reguli de acces, stocare a dovezilor și posibilitatea de administrare a incidentelor. Tocmai această legătură dintre componente transformă modelul de inteligență artificială într-o aplicație utilizabilă.

Concret, domeniul impune șase cerințe care se regăsesc direct în proiectarea sistemului. Detecția vizuală automată este necesară deoarece monitorizarea manuală continuă este dificilă și consumatoare de timp; în lucrare, ea este acoperită de detectorul de deșeuri. Analiza temporală este necesară deoarece un obiect detectat într-un singur cadru nu confirmă un incident; ea este realizată prin corelarea apariției obiectului cu prezența persoanei și cu persistența obiectului în scenă. Validarea umană rămâne obligatorie deoarece scenele reale pot produce detecții greșite sau interpretări ambigue, astfel încât fiecare incident este confirmat, respins sau arhivat de utilizatorul responsabil. Dovezile locale permit verificarea ulterioară a evenimentului semnalat, prin imagine reprezentativă, clip scurt și metadate. Controlul accesului separă funcționalitățile pe roluri, deoarece nu toate acțiunile trebuie să fie disponibile tuturor utilizatorilor. În sfârșit, protecția datelor impune păstrarea controlată a dovezilor, deoarece fluxurile video pot conține persoane și informații sensibile.

\section{Sisteme inteligente pentru detecția obiectelor abandonate}

Sistemele inteligente bazate pe viziune artificială urmăresc extragerea automată a informațiilor relevante din imagini sau secvențe video. În cazul detecției obiectelor, scopul este localizarea și clasificarea unor entități vizibile în cadru. Modelele moderne de detecție folosesc rețele neuronale convoluționale, care au devenit o soluție standard în analiza imaginilor datorită capacității lor de a învăța caracteristici vizuale direct din date \cite{lecun1998gradient}, \cite{krizhevsky2012imagenet}.

În monitorizarea video, unde cadrele trebuie analizate succesiv, echilibrul dintre acuratețe și timp de procesare devine esențial. Din acest motiv, lucrarea utilizează familia YOLO, orientată către detecție rapidă, în implementarea Ultralytics YOLOv8, care oferă un cadru accesibil pentru antrenare, evaluare și inferență \cite{redmon2016yolo}, \cite{ultralytics2023yolov8}. Alegerea este importantă nu doar pentru performanța experimentală, ci și pentru integrarea într-o aplicație locală, unde resursele hardware pot fi limitate; clasificarea detaliată a familiilor de modele de detecție și justificarea teoretică a acestei alegeri sunt prezentate în capitolul următor.

Cercetări recente propun și variante optimizate ale familiei YOLO pentru detectarea deșeurilor în scene complexe, inclusiv în zone urbane sau pe suprafețe cu fundal variabil \cite{liu2024ecodetect}. Aceste rezultate susțin direcția aleasă în lucrare: folosirea unui detector rapid, urmată de verificare umană și de reguli suplimentare de interpretare temporală. Într-o aplicație de monitorizare, performanța modelului trebuie analizată împreună cu stabilitatea fluxului, calitatea dovezilor salvate și claritatea interfeței de validare.

Un sistem pentru detectarea obiectelor abandonate nu poate depinde însă doar de detectorul de deșeuri. Același obiect poate fi detectat în mai multe cadre consecutive, poate dispărea temporar din cauza unei persoane care trece prin fața camerei sau poate fi confundat cu elemente ale fundalului. Pentru a interpreta corect un posibil incident, sistemul trebuie să urmărească evoluția scenei. Algoritmii de tracking multi-obiect, cum este ByteTrack, evidențiază importanța asocierii detecțiilor între cadre pentru înțelegerea mișcării și persistenței obiectelor \cite{zhang2022bytetrack}.

În această lucrare, ideea de obiect abandonat este modelată printr-o succesiune de observații. O persoană este prezentă în zona monitorizată, un obiect de tip deșeu apare sau devine stabil în cadru, iar scena indică posibilitatea ca obiectul să fi fost lăsat în urmă. Această logică nu pretinde să reconstruiască perfect intenția persoanei. Scopul ei este mai practic: reducerea numărului de momente pe care persoana responsabilă trebuie să le verifice manual și păstrarea dovezilor relevante pentru fiecare alertă.

Clasificarea materialului completează detecția. Pentru utilizatorul responsabil, informația că un obiect este estimat ca plastic, hârtie, sticlă sau metal poate ajuta la descrierea incidentului și la organizarea dovezilor. În sistemul propus, clasificarea este aplicată după localizarea obiectului, asupra regiunii de imagine corespunzătoare detecției. Astfel, incidentul nu conține doar poziția obiectului, ci și o descriere mai utilă pentru verificare.

\section{Concepte și termeni utilizați în lucrare}

Pentru claritatea lucrării, sunt definite principalele concepte folosite în capitolele următoare. Prin cadru video se înțelege o imagine individuală extrasă dintr-un flux video. Prin detecție se înțelege identificarea automată a unui obiect într-un cadru, împreună cu poziția acestuia. Poziția este reprezentată printr-o casetă de delimitare, numită frecvent bounding box, care indică zona din imagine în care modelul estimează prezența obiectului.

Fiecare detecție este însoțită de un scor de încredere. Acest scor nu reprezintă o certitudine absolută, ci o estimare numerică produsă de model. Un prag de încredere prea mic poate produce multe alerte false, iar un prag prea mare poate omite obiecte reale. Din acest motiv, pragurile de detecție trebuie alese în funcție de scenariu și validate experimental. În lucrare, aceste aspecte sunt tratate în capitolele dedicate antrenării și testării modelelor.

Prin clasificare se înțelege atribuirea unei etichete unui obiect sau unei imagini. În cazul sistemului propus, clasificarea materialului se aplică obiectului detectat și urmărește categorii precum plastic, hârtie, sticlă, metal sau alte materiale. Prin tracking se înțelege urmărirea unei entități de la un cadru la altul, pentru a observa dacă aceasta persistă, se deplasează sau dispare. Trackingul este important deoarece incidentul nu poate fi interpretat corect dintr-o singură imagine.

Termenul de incident desemnează o situație semnalată de sistem în care apariția unui obiect de tip deșeu, corelată cu contextul temporal, sugerează o posibilă acțiune de aruncare sau abandonare. Incidentul nu este echivalent cu o constatare oficială. El este o alertă tehnică, însoțită de dovezi locale, care trebuie analizată de o persoană responsabilă. Prin dovadă locală se înțelege ansamblul format din metadate, imagine reprezentativă și clip video scurt asociat evenimentului.

În procesul de validare pot apărea două tipuri de erori. Un fals pozitiv apare atunci când sistemul semnalează un incident care, după verificare, nu reprezintă o aruncare ilegală reală. Un fals negativ apare atunci când un incident real nu este semnalat de sistem. Ambele tipuri de erori sunt importante. Falsurile pozitive afectează încrederea utilizatorului și pot încărca inutil panoul de administrare, iar falsurile negative reduc utilitatea practică a sistemului. De aceea, evaluarea trebuie să includă atât metrici ale modelelor, cât și observații asupra comportamentului sistemului pe clipuri video.

Intrarea principală a sistemului este un flux video sau un fișier video încărcat de utilizator. Ieșirea practică este o listă de incidente care pot fi vizualizate, filtrate și validate într-o interfață web locală. Fiecare incident conține informații despre momentul apariției, materialul estimat, scorurile relevante, statusul de validare și fișierele media asociate. Această structură permite legarea componentei de inteligență artificială de o aplicație utilizabilă, în care rezultatele nu rămân simple predicții, ci devin elemente gestionabile într-un flux de lucru.

Pe lângă acești termeni de bază, capitolele următoare folosesc câteva noțiuni specifice algoritmului temporal. Prin zona persoanei se înțelege ultima regiune din cadru în care o persoană a fost observată înainte de a părăsi scena; această zonă, extinsă cu o marjă procentuală, delimitează spațiul în care apariția unui obiect nou este considerată relevantă. Prin fereastra de monitorizare se înțelege intervalul de timp în care zona persoanei este supravegheată după plecarea acesteia. Fereastra de confirmare desemnează intervalul scurt dintre identificarea unui candidat de incident și emiterea efectivă a alertei, folosit pentru a anula cazurile în care persoana revine și ridică obiectul. Perioada de cooldown reprezintă intervalul de după o alertă în care alertele suplimentare sunt suprimate, pentru a evita raportarea aceluiași eveniment de mai multe ori.

În concluzie, domeniul abordat de lucrare se află la intersecția dintre monitorizarea spațiilor publice, viziunea artificială și dezvoltarea aplicațiilor web. Capitolul de față a delimitat problema, a prezentat necesitatea monitorizării automate și a introdus termenii utilizați în continuare. Capitolul următor va detalia fundamentele teoretice și tehnologiile care permit implementarea soluției propuse.
